% GENERATED FILE. Edit canonical lessons, then run npm run generate:books.
% canonical-source-hash: fnv1a64:1eba9d993ebbb860
% canonical-lessons: GU-C172-phiku, GU-C172-vasi, GU-C172-balvu, GU-C172-narangi, GU-C172-ras

\chapter{Stale, burnt and fresh juice}
\label{ch:gu-stale-burnt-and-fresh-juice}
\hlchaptermodality{172}

\begin{chapteropening}
\textbf{By the end of this chapter:} \emph{I can say that food is bland, stale or burnt.}

Complete the last lesson of chapter 172: I can say that food is bland, stale or burnt.
\end{chapteropening}

\section[\texorpdfstring{\gu{ફીકું} (phīkũ)}{ફીકું (phīkũ)}]{\gu{ફીકું} (phīkũ) — bland, pale}

\label{lesson:GU-C172-phiku}



\begin{quote}
Before the new one: say the Gujarati for salty, then the Gujarati for sweet.
\end{quote}

\subsection*{You'll want to know: \gu{ફીકું}}
\textbf{\gu{ફીકું}} — \emph{phīkũ} — \textquotedblleft{}bland, pale\textquotedblright{}.

\begin{quote}
\textbf{\gu{આ} \gu{શાક} \gu{ફીકું} \gu{છે}.} — \emph{ā śāk phīkũ che} — \textquotedblleft{}This curry is bland.\textquotedblright{}
\end{quote}

\begin{cousinweb}
\textbf{\gu{ફીકું}} (\emph{phīkũ}) is a tadbhava. It comes from \emph{*phikka}, \textquotedblleft{}defective\textquotedblright{}, a reconstructed form, and CDIAL glosses the Gujarati word \textquotedblleft{}pale, tasteless\textquotedblright{} (CDIAL 9037).
\end{cousinweb}

\begin{grammarlens}[title={What you've built}]
One more word for everyday life at A2.
\end{grammarlens}

\subsection*{Your turn}
\begin{itemize}
\raggedright
  \item \emph{Say it:} \emph{phīkũ}
  \item \emph{Say it:} \emph{phīkũ}, once more
  \item \emph{Say it:} \emph{ā śāk phīkũ che} — This curry is bland.
  \item \emph{Recall:} say the Gujarati for salty, then the Gujarati for sweet, then say \emph{phīkũ} again
\end{itemize}

\begin{tcolorbox}[breakable,colback=teal!4,colframe=teal!35!black,title={Before you move on}]
What does \gu{ફીકું} mean? (\textquotedblleft{}Bland, pale\textquotedblright{}.) Say it once more.
\end{tcolorbox}
\section[\texorpdfstring{\gu{વાસી} (vāsī)}{વાસી (vāsī)}]{\gu{વાસી} (vāsī) — stale}

\label{lesson:GU-C172-vasi}



\begin{quote}
Before the new one: say the Gujarati for sweet, then the Gujarati for bland.
\end{quote}

\subsection*{You'll want to know: \gu{વાસી}}
\textbf{\gu{વાસી}} — \emph{vāsī} — \textquotedblleft{}stale\textquotedblright{}.

\begin{quote}
\textbf{\gu{આ} \gu{રોટલી} \gu{વાસી} \gu{છે}.} — \emph{ā roṭlī vāsī che} — \textquotedblleft{}This rotli is stale.\textquotedblright{}
\end{quote}

\begin{cousinweb}
\textbf{\gu{વાસી}} (\emph{vāsī}) is a tadbhava. CDIAL lists it, glossed \textquotedblleft{}stale\textquotedblright{}, under Sanskrit \emph{vāsita}, \textquotedblleft{}infused, scented\textquotedblright{} (CDIAL 11604).
\end{cousinweb}

\begin{grammarlens}[title={What you've built}]
One more word for everyday life at A2.
\end{grammarlens}

\subsection*{Your turn}
\begin{itemize}
\raggedright
  \item \emph{Say it:} \emph{vāsī}
  \item \emph{Say it:} \emph{vāsī}, once more
  \item \emph{Say it:} \emph{ā roṭlī vāsī che} — This rotli is stale.
  \item \emph{Recall:} say the Gujarati for sweet, then the Gujarati for bland, then say \emph{vāsī} again
\end{itemize}

\begin{tcolorbox}[breakable,colback=teal!4,colframe=teal!35!black,title={Before you move on}]
What does \gu{વાસી} mean? (\textquotedblleft{}Stale\textquotedblright{}.) Say it once more.
\end{tcolorbox}
\section[\texorpdfstring{\gu{બળવું} (baḷvũ)}{બળવું (baḷvũ)}]{\gu{બળવું} (baḷvũ) — to get burnt, to burn}

\label{lesson:GU-C172-balvu}



\begin{quote}
Before the new one: say the Gujarati for bland, then the Gujarati for stale.
\end{quote}

\subsection*{You'll want to know: \gu{બળવું}}
\textbf{\gu{બળવું}} — \emph{baḷvũ} — \textquotedblleft{}to get burnt, to burn\textquotedblright{}.

\begin{quote}
\textbf{\gu{રોટલી} \gu{તવા} \gu{ઉપર} \gu{બળે} \gu{છે}.} — \emph{roṭlī tavā upar baḷe che} — \textquotedblleft{}The rotli burns on the griddle.\textquotedblright{}
\end{quote}

\begin{cousinweb}
\textbf{\gu{બળવું}} (\emph{baḷvũ}) is a tadbhava. It comes from \emph{*dvalati}, \textquotedblleft{}burns\textquotedblright{}, a reconstructed form beside Sanskrit \emph{jvalati}, and CDIAL lists the Gujarati \emph{baḷvũ}, \textquotedblleft{}burns\textquotedblright{} (CDIAL 6654).
\end{cousinweb}

\begin{grammarlens}[title={What you've built}]
One more word for everyday life at A2.
\end{grammarlens}

\subsection*{Your turn}
\begin{itemize}
\raggedright
  \item \emph{Say it:} \emph{baḷvũ}
  \item \emph{Say it:} \emph{baḷvũ}, once more
  \item \emph{Say it:} \emph{roṭlī tavā upar baḷe che} — The rotli burns on the griddle.
  \item \emph{Recall:} say the Gujarati for bland, then the Gujarati for stale, then say \emph{baḷvũ} again
\end{itemize}

\begin{tcolorbox}[breakable,colback=teal!4,colframe=teal!35!black,title={Before you move on}]
What does \gu{બળવું} mean? (\textquotedblleft{}To get burnt, to burn\textquotedblright{}.) Say it once more.
\end{tcolorbox}
\section[\texorpdfstring{\gu{નારંગી} (nārangī)}{નારંગી (nārangī)}]{\gu{નારંગી} (nārangī) — orange}

\label{lesson:GU-C172-narangi}



\begin{quote}
Before the new one: say the Gujarati for stale, then the Gujarati for to get burnt.
\end{quote}

\subsection*{You'll want to know: \gu{નારંગી}}
\textbf{\gu{નારંગી}} — \emph{nārangī} — \textquotedblleft{}orange\textquotedblright{}.

\begin{quote}
\textbf{\gu{નારંગી} \gu{ખાટી} \gu{અને} \gu{મીઠી} \gu{છે}.} — \emph{nārangī khāṭī ane mīṭhī che} — \textquotedblleft{}The orange is sour and sweet.\textquotedblright{}
\end{quote}

\begin{cousinweb}
\textbf{\gu{નારંગી}} (\emph{nārangī}) is a tadbhava. It comes from Sanskrit \emph{nāraṅga}, \emph{nāraṅgī}, \textquotedblleft{}orange tree\textquotedblright{}, and CDIAL glosses the Gujarati word \textquotedblleft{}the tree and fruit\textquotedblright{} (CDIAL 7073).
\end{cousinweb}

\begin{grammarlens}[title={What you've built}]
One more word for everyday life at A2.
\end{grammarlens}

\subsection*{Your turn}
\begin{itemize}
\raggedright
  \item \emph{Say it:} \emph{nārangī}
  \item \emph{Say it:} \emph{nārangī}, once more
  \item \emph{Say it:} \emph{nārangī khāṭī ane mīṭhī che} — The orange is sour and sweet.
  \item \emph{Recall:} say the Gujarati for stale, then the Gujarati for to get burnt, then say \emph{nārangī} again
\end{itemize}

\begin{tcolorbox}[breakable,colback=teal!4,colframe=teal!35!black,title={Before you move on}]
What does \gu{નારંગી} mean? (\textquotedblleft{}Orange\textquotedblright{}.) Say it once more.
\end{tcolorbox}
\section[\texorpdfstring{\gu{રસ} (ras)}{રસ (ras)}]{\gu{રસ} (ras) — juice}

\label{lesson:GU-C172-ras}



\begin{quote}
Before the new one: say the Gujarati for to get burnt, then the Gujarati for orange.
\end{quote}

\subsection*{You'll want to know: \gu{રસ}}
\textbf{\gu{રસ}} — \emph{ras} — \textquotedblleft{}juice\textquotedblright{}.

\begin{quote}
\textbf{\gu{હું} \gu{સવારે} \gu{નારંગીનો} \gu{રસ} \gu{પીઉં} \gu{છું}.} — \emph{hũ savāre nārangīno ras pīũ chũ} — \textquotedblleft{}I drink orange juice in the morning.\textquotedblright{}
\end{quote}

\begin{cousinweb}
\textbf{\gu{રસ}} (\emph{ras}) is inherited from Sanskrit \emph{rasa}, \textquotedblleft{}sap or juice of plants; taste\textquotedblright{}, with only its final vowel lost, and CDIAL glosses the Gujarati word \textquotedblleft{}sap, juice\textquotedblright{} (CDIAL 10650).
\end{cousinweb}

\begin{grammarlens}[title={What you've built}]
One more word for everyday life at A2.
\end{grammarlens}

\subsection*{Your turn}
\begin{itemize}
\raggedright
  \item \emph{Say it:} \emph{ras}
  \item \emph{Say it:} \emph{ras}, once more
  \item \emph{Say it:} \emph{hũ savāre nārangīno ras pīũ chũ} — I drink orange juice in the morning.
  \item \emph{Recall:} say the Gujarati for to get burnt, then the Gujarati for orange, then say \emph{ras} again
\end{itemize}

\begin{tcolorbox}[breakable,colback=teal!4,colframe=teal!35!black,title={Before you move on}]
What does \gu{રસ} mean? (\textquotedblleft{}Juice\textquotedblright{}.) Say it once more.
\end{tcolorbox}
